\section{Custom Compensating Control: MORI}
\label{sec:compensating-control}

Chapter~6 demonstrated that the integrity monitor could recognise a changed
privileged target without identifying the proof-of-concept sample. That result
provided visibility, but it occurred after the protected file view had already
changed and did not prevent the resulting UID~0 transition. RQ4 therefore moves
the defensive decision earlier in the path: can a temporary control interrupt
the known exploitation sequence while the vulnerable kernel is still in use,
and can it avoid removing benign functionality that is not required to be
denied?

The experiment approached this question in two stages. MORI Guard v1 applied a
broad module-loading restriction to \code{algif\_aead}. It stopped the tested
proof of concept on kernel \code{6.8.0-116-generic}, but a benign unprivileged
AF\_ALG AEAD bind failed for the same reason. The result was useful as a safe
baseline: it showed that removing access to the complete interface interrupted
the path, while also making the compatibility cost directly visible.

MORI v2 was then developed as a more selective eBPF and Linux Security Module
(LSM) control. Rather than classifying any one syscall as malicious, it combined
AF\_ALG AEAD activity, a root-owned SUID-file-to-pipe splice, and the associated
privileged-file read in a shared process and thread context. The implementation
was introduced through observation-only checkpoints, negative controls, a
shadow policy, active denial, structured telemetry, adversarial regression
testing, and finally a lifecycle correction in MORI v2.7.

The final control returned \code{EPERM} at the privileged-file permission check
inside the correlated splice operation. The preserved final validation shows
the frozen v2.7 executable denying both the laboratory Python proof of concept
and an independently implemented C proof of concept. In both tests,
\code{labuser} remained UID~1001 and the SHA-256 digest of \code{/usr/bin/su}
remained equal to its known-good value. The result is deliberately narrower
than a claim of universal prevention: MORI v2.7 enforces an experimentally
derived policy for the tested same-process Copy Fail paths and remains a
compensating control rather than a correction to the kernel defect.

This chapter documents the control objective, evidential boundary, broad v1
baseline, incremental v2 development, final enforcement architecture,
lifecycle regression, compatibility controls, adversarial validation, and
remaining limitations. It concludes by answering RQ4 and establishing the
boundary for the vendor comparison in Chapter~8.

\subsection{Control Objective, Success Criteria, and Evidence Boundary}

The compensating-control objective was not merely to make the public Python
script terminate with an error. A useful result had to show where the
interruption occurred, preserve the unprivileged identity, protect the selected
SUID target, and distinguish selective prevention from wholesale interface
removal. The evaluation criteria are summarised in
\cref{tab:mori-success-criteria}.

\begin{table}[H]
    \centering
    \caption{Success criteria used for the custom compensating control.}
    \label{tab:mori-success-criteria}
    \small
    \begin{tabularx}{\textwidth}{@{}>{\bfseries\raggedright\arraybackslash}p{0.24\textwidth}>{\raggedright\arraybackslash}p{0.30\textwidth}Y@{}}
        \toprule
        Criterion & Required observation & Evidential meaning \\
        \midrule
        Known path interrupted &
        The tested proof of concept receives a denial before obtaining a
        privileged shell. &
        Demonstrates prevention for that execution path, not correction of the
        underlying vulnerability. \\

        Privilege boundary preserved &
        Post-test identity remains \code{uid=1001(labuser)}. &
        Excludes an apparent error followed by a successful UID~0 transition. \\

        Privileged target preserved &
        \code{/usr/bin/su} retains the known-good mounted-view SHA-256 digest. &
        Shows that the selected target did not enter the changed state observed
        in Chapter~5. \\

        Benign interface use retained &
        An unprivileged AF\_ALG AEAD control can still reach the tested
        socket, bind, or accept operation. &
        Distinguishes selective enforcement from disabling
        \code{algif\_aead}; it does not prove a complete cryptographic
        round trip. \\

        Individual signals remain permitted &
        AEAD-only, splice-only, and deliberately separated process activity do
        not reach the deny condition. &
        Tests whether the control enforces the correlation rather than treating
        every prerequisite as malicious. \\

        Reproduced defect is regression-tested &
        A test that bypasses an earlier policy version is retained and repeated
        against the corrected candidate. &
        Demonstrates that the final change addresses the observed lifecycle
        failure rather than only changing labels or telemetry. \\
        \bottomrule
    \end{tabularx}
\end{table}

The Phase~04 package contains several evidence classes. MORI Guard v1 is
preserved through machine-readable configuration, blocker source, deployment
state, compatibility output, final retest output, journal events, and
screenshots. MORI v2 preserves complete implementation checkpoints from
v2.2.0 through v2.6.2, screenshots for the earlier v2.1.x work, final v2.7
source and build objects, and the regression and proof-of-concept captures.
This progression is important because the failed designs and compatibility
results are part of the control evaluation rather than discarded development
noise.

The current MORI v2 public manifest contains 108 SHA-256 entries. It includes
the historical checkpoints, the v2.7 source and build set, the current
provenance records, and screenshots~00--57, with the two \code{37a} and
\code{37b} captures retained separately. The older inner
\code{mori-v2/README.md} still reports an 85-object inventory from the earlier
v2.6.2 package. It is therefore used for historical design notes rather than as
the final inventory. The manifest and the current Phase~04 index are the
authoritative package boundaries used in this chapter.

The frozen v2.7 executable is identified by:

\begin{lstlisting}
7fe7b609b90e164f3ae4a9c025363399a8a63bf2a4338db2ecc62cfd2682d039
\end{lstlisting}

That digest is visible before the final Python validation and is repeated in
the frozen artifact manifest and later provenance record. The associated
source files are preserved with the following digests:

\begin{lstlisting}
05e0a387ce084d2aef121351c9af7907f85696307d9f08ccdf6d9191f91a035d
    source/mori_observer.bpf.c
2ee01194b488e757dd4c1062a4123aec8aac5c5cbd74767ebc3a938f886939e6
    source/mori_observer.c
\end{lstlisting}

One chronology boundary must be kept explicit. The final artifact-provenance
record was captured later, when the host reported kernel
\code{6.8.0-137-generic}. It verifies the hashes of the transferred frozen
source and build set; it is not used as evidence of the earlier mitigation-test
kernel. The v1 machine-readable retest and v2.6.2 pre-test capture directly
record \code{6.8.0-116-generic}. The final v2.7 captures identify the tested
binary by hash and show the denial and post-test state, but do not repeat
\code{uname}. The project chronology places those captures in Phase~04 before
the separate vendor-remediation validation. This distinction prevents the
later provenance environment from being substituted for the actual test
sequence.

The evidence supports conclusions about the two tested proof-of-concept
implementations and the stated regression controls. It does not establish a
population-wide prevention rate, production compatibility, or protection
against every possible variation of the vulnerability.

\subsection{MORI Guard v1: Broad Interruption and Compatibility Cost}

MORI Guard v1 used the normal \code{modprobe} configuration mechanism to
redirect requests for \code{algif\_aead} to a root-owned blocker script:

\begin{lstlisting}
install algif_aead /usr/local/sbin/mori-block-algif-aead
\end{lstlisting}

The blocker wrote a structured journal entry, printed the presentation message
to the requesting process when possible, and exited with status~1. The
deployment record shows that \code{modprobe -n -v algif\_aead} resolved to the
custom command, that \code{algif\_aead} was not loaded, and that the existing
MORI integrity monitor remained active. This was an administrative loading
policy, not a change to the vulnerable kernel code.

The final v1 retest was performed by \code{labuser} on
\code{6.8.0-116-generic}. Before the attempt, the account was UID~1001 and
\code{/usr/bin/su} had its expected digest:

\begin{lstlisting}
c74311fe5636b7d7f9a56239fa8adeeab12ba86fe7d41b91afa85bf9bbdae78b
\end{lstlisting}

The proof of concept terminated with \code{FileNotFoundError}, returned exit
status~1, and did not produce a privileged shell. The post-attempt identity
remained \code{labuser}, the target digest was unchanged, and
\code{algif\_aead} remained blocked. The journal recorded the blocker decision,
while the integrity-monitor query returned no event because the protected
target did not change. Within the normal module-loading path exercised by the
test, v1 therefore interrupted the known exploit prerequisite before the
privilege boundary was crossed.

The same mechanism also exposed its compatibility cost. A separate unprivileged
control attempted to bind an AF\_ALG AEAD socket to \code{gcm(aes)}. With v1
active, it failed with \code{FileNotFoundError}, \code{errno=2}, and
\code{No such file or directory}. The failure was expected, but it established
that the control could not distinguish the Copy Fail sequence from a benign
consumer requesting the same kernel interface.

\begin{table}[H]
    \centering
    \caption{Security and compatibility result of MORI Guard v1.}
    \label{tab:mori-v1-result}
    \small
    \begin{tabularx}{\textwidth}{@{}>{\bfseries\raggedright\arraybackslash}p{0.27\textwidth}>{\raggedright\arraybackslash}p{0.25\textwidth}Y@{}}
        \toprule
        Test & Direct result & Interpretation \\
        \midrule
        Python Copy Fail retest &
        PoC failed; UID remained 1001; target hash unchanged. &
        The tested path was interrupted while the vulnerable kernel was
        retained. \\

        Benign \code{gcm(aes)} AEAD bind &
        Failed while the guard was active. &
        The protection was obtained by removing the tested interface from both
        malicious and benign callers. \\

        MORI integrity monitor &
        No target-mismatch event. &
        Consistent with prevention before the monitored target changed; not an
        independent proof of why the PoC failed. \\
        \bottomrule
    \end{tabularx}
\end{table}

The result answers only the first half of RQ4. MORI Guard v1 prevented the
known execution path, but it did so by imposing the exact broad functionality
loss that the second half of the question asks the project to avoid. It
therefore became the compatibility baseline and the motivation for MORI v2.

The scope of the v1 mechanism is also narrower than the phrase ``module
disabled'' may imply. The test validated normal \code{modprobe} resolution
while the module was unloaded. It did not show prevention of direct insertion
by an already privileged administrator, protection when the functionality is
built into the kernel, or resistance to a root user changing the configuration
and blocker script. Those are acceptable boundaries for the temporary
laboratory control, but they prevent it from being described as a kernel fix.

\subsection{From a False Positive to a Selective Policy}

MORI v2 did not begin as an enforcement engine. It was developed by first
observing each candidate signal, then testing whether state could be shared
between hooks, and only later converting the correlation into a denial. The
development sequence is summarised in
\cref{tab:mori-development-progression}.

\begin{table}[H]
    \centering
    \caption{MORI v2 development and the question answered at each checkpoint.}
    \label{tab:mori-development-progression}
    \scriptsize
    \begin{tabularx}{\textwidth}{@{}>{\bfseries}p{0.10\textwidth}>{\raggedright\arraybackslash}p{0.28\textwidth}Y@{}}
        \toprule
        Version & Change & Experimental lesson \\
        \midrule
        v2.1 & Observe non-root reads of root-owned SUID files through
        \code{lsm/file\_permission}. &
        MORI Monitor's legitimate periodic hashing immediately triggered the
        observer; a SUID read alone was too broad. \\

        v2.1.1 & Add a \code{trusted\_tgids} map populated from the monitor's
        service cgroup. &
        The known monitor could be suppressed, while an untrusted SUID read
        remained visible. Cgroup membership alone was still an incomplete
        trust decision. \\

        v2.1.2 & Authenticate the monitor path, file type, owner, permissions,
        SHA-256 digest, cgroup membership, and exact argv entry. &
        Tampering caused trust to be refused. MORI still attached with an empty
        exemption map, making the monitor observable rather than silently
        trusting it. \\

        v2.2 & Observe AF\_ALG \code{aead} bind attempts through
        \code{lsm/socket\_bind}. &
        The hook could see the candidate signal. With v1 disabled, an
        unprivileged \code{gcm(aes)} bind also succeeded, demonstrating that
        observation need not remove the interface. \\

        v2.3 & Record recent AEAD state by TGID and correlate it with a SUID
        read over an experimental ten-second interval. &
        Same-process temporal correlation worked and expired as designed, but
        time proximity did not establish malicious intent. \\

        v2.4 & Observe root-owned SUID-file-to-pipe operations at
        \code{fentry/do\_splice} and share state with the LSM program. &
        Cross-hook state was visible during the same operation; AEAD, splice,
        and SUID-read signals could be combined for one TGID. \\

        v2.5 & Emit \code{WOULD-DENY-}\allowbreak
        \code{AEAD-SPLICE-SUID} while preserving the original LSM return value. &
        The combined path reached the shadow decision, while AEAD-only and
        splice-only controls remained allowed. \\

        v2.6 & Return \code{-EPERM} for the validated correlation. &
        The benign AEAD bind succeeded, but the correlated SUID-to-pipe splice
        was denied with \code{errno=1}. \\

        v2.6.1 & Add a BPF ring buffer containing timestamp, TGID, UID, and
        event type. &
        Userspace received structured deny telemetry, while enforcement
        remained independent of successful event delivery. \\

        v2.6.2 & Add best-effort interactive-terminal notification and extend
        adversarial and target regression testing. &
        The presentation path was separated from the LSM decision. Later
        testing also exposed a same-TGID expiry weakness in the timed state. \\

        v2.7 & Replace the enforcement timer with socket-lifetime AEAD state and
        exact-invocation splice state. &
        The retained 12-second regression was denied, followed by final Python
        and independent C proof-of-concept validation. \\
        \bottomrule
    \end{tabularx}
\end{table}

The first false positive was valuable because it prevented the SUID-read signal
from becoming an unquestioned enforcement rule. The integrity monitor performed
exactly the same low-level action that the prototype observed: it read
root-owned SUID executables. Suppressing it by a hard-coded PID, process name,
or UID would have created a weak exemption. The final userspace loader instead
required a root-owned regular monitor file, no group or world write bit, the
expected SHA-256 digest, an exact command-line argument, and membership in the
dedicated service cgroup before populating the trusted map.

This trust mechanism is an exemption, not the source of the security decision.
If verification fails, MORI emits a warning, attaches its BPF programs, and
leaves the trust map empty. That behaviour favours continued observation over a
silent exemption. It can produce additional monitor noise, but it does not
disable the compensating control simply because the companion monitor could not
be authenticated.

The shadow-policy stage provided a second safety boundary. The combined rule
was allowed to reach a \code{WOULD-DENY} state before any operation was
rejected. AEAD-only and SUID-splice-only negative controls did not reach that
state, whereas the deliberately combined sequence did. Only after this policy
shape was visible did v2.6 change the LSM return value to \code{-EPERM}.
Consequently, active denial was based on a previously observed correlation
rather than introduced at the same time as the first multi-hook test.

\subsection{MORI v2.7 Enforcement Architecture}

The frozen v2.7 candidate consists of a userspace loader and a CO-RE BPF object.
The loader opens and loads the object, authenticates the optional MORI Monitor
exemption, attaches the BPF programs, and polls a ring buffer for deny events.
The kernel-side programs make the enforcement decision. The four principal
state maps are shown in \cref{tab:mori-v27-state}.

\begin{table}[H]
    \centering
    \caption{Kernel-side state used by MORI v2.7.}
    \label{tab:mori-v27-state}
    \small
    \begin{tabularx}{\textwidth}{@{}>{\bfseries\raggedright\arraybackslash}p{0.27\textwidth}>{\raggedright\arraybackslash}p{0.25\textwidth}Y@{}}
        \toprule
        Map & Key and value & Purpose \\
        \midrule
        \code{trusted\_tgids} & TGID $\rightarrow$ one-byte trust flag &
        Suppresses reads from the authenticated MORI integrity-monitor process. \\

        \code{aead\_socket\_state} & BPF socket-local storage $\rightarrow$
        owner TGID, state, and counted flag &
        Associates an observed AF\_ALG AEAD socket with the process that bound
        it and follows the tracked socket until release. \\

        \code{active\_aead\_tgids} & TGID $\rightarrow$ active-socket count &
        Allows enforcement to ask whether the current process owns at least one
        live tracked AEAD socket. \\

        \code{splice\_active\_tids} & TID $\rightarrow$ one-byte active flag &
        Exists only while that exact thread is inside a qualifying
        root-SUID-file-to-pipe \code{do\_splice} invocation. \\

        \code{events} & Ring-buffer records &
        Carries timestamp, TGID, UID, and deny type to userspace on a
        best-effort basis. \\
        \bottomrule
    \end{tabularx}
\end{table}

At \code{lsm/socket\_bind}, MORI first preserves any non-zero decision returned
by a previous LSM program. It then checks for address family~38 (AF\_ALG) and
an algorithm type whose fixed field begins with \code{aead}. The tracked socket
is associated with the current TGID, and the process's active-AEAD count is
incremented. A separate \code{lsm/socket\_accept} program reports an accept only
when the parent socket already carries MORI state. At
\code{fentry/\_\_sock\_release}, the active count is decremented and the
socket-local state is destroyed with the socket.

The implementation therefore makes the lifetime of the tracked socket, rather
than a fixed number of seconds, the relevant AEAD state. The accept event is
useful telemetry, but the frozen source activates and counts the state in the
bind hook. This distinction matters because \code{socket\_bind} observes the
attempt before userspace independently confirms successful completion; it is
retained as a limitation later in this chapter.

At \code{fentry/do\_splice}, the program checks that the caller is non-root,
the input inode is root-owned and carries the owner SUID bit, and the output
inode is a FIFO. If those conditions hold, the current TID is inserted into
\code{splice\_active\_tids}. The corresponding \code{fexit/do\_splice}
program deletes that entry. The state is thus scoped to the exact thread and
the exact invocation instead of remaining available for a grace period after
the operation.

The decision point is \code{lsm/file\_permission}. For a non-root read of a
root-owned SUID file, the program first checks the trusted-monitor exemption.
It then looks for both of the following:

\begin{lstlisting}
active_aead_tgids[current TGID] > 0
splice_active_tids[current TID] == 1
\end{lstlisting}

When both are present, MORI submits a ring-buffer event if space is available,
writes a trace message, and returns \code{-1}, corresponding to
\code{-EPERM}. If either condition is absent, the previous LSM result is
preserved. The policy interrupts the file read used inside the qualifying
splice; it does not wait for the later page-cache change or privileged-program
execution observed in Chapters~5 and~6.

The tested policy boundary is summarised in
\cref{tab:mori-v27-policy-boundary}.

\begin{table}[H]
    \centering
    \caption{Decision boundary implemented by the frozen MORI v2.7 source.}
    \label{tab:mori-v27-policy-boundary}
    \small
    \begin{tabularx}{\textwidth}{@{}>{\raggedright\arraybackslash}p{0.43\textwidth}>{\bfseries\raggedright\arraybackslash}p{0.16\textwidth}Y@{}}
        \toprule
        Observed condition & Decision & Reason \\
        \midrule
        Ordinary socket activity or non-AEAD AF\_ALG activity & Allowed &
        No tracked AEAD socket is created by the tested predicate. \\

        AF\_ALG AEAD activity without a qualifying splice & Allowed &
        The splice-invocation state is absent. \\

        Root-owned SUID file spliced to a pipe without same-process active AEAD
        state & Allowed &
        The active-AEAD count is absent or zero. \\

        AEAD activity in process A and SUID-to-pipe splice in process B &
        Allowed &
        AEAD state is isolated by TGID rather than correlated across processes. \\

        Live tracked AEAD socket in the same TGID, with the current TID inside
        the qualifying splice and reading the SUID input & Denied with EPERM &
        Both enforcement states are present at the protected read. \\

        Authenticated MORI Monitor read & Allowed &
        The trusted TGID exemption is evaluated before correlation. \\

        Root caller & Allowed by this policy &
        The control is designed to interrupt the pre-escalation non-root path,
        not constrain an already privileged process. \\
        \bottomrule
    \end{tabularx}
\end{table}

Structured telemetry is deliberately downstream of the decision. If the ring
buffer is full or the userspace process is unavailable, reservation can fail,
but the LSM program still returns \code{EPERM}. When userspace does receive the
event, it prints the TGID and UID and optionally searches
\code{/proc/<tgid>/fd/} for descriptors~2, 1, and~0. Each candidate is opened
independently and checked with \code{isatty()} before the cosmetic MORI message
is written. Failure to find or write a terminal changes neither the kernel
decision nor the denied process's return value.

\subsection{Regression-Driven Lifecycle Correction}

MORI v2.6.2 used timestamps keyed by TGID. AEAD state and splice state were
considered relevant only within an experimental ten-second interval. That
design passed the original short positive tests, but later adversarial testing
kept the same process and AEAD socket alive, waited 12~seconds, and then spliced
\code{/usr/bin/fusermount3} into a pipe. The trace showed the expected AEAD and
splice observations, yet userspace reported:

\begin{lstlisting}
same-TGID expiry test: tgid=8358
AEAD armed; waiting 12 seconds...
splice allowed after expiry: 32 bytes
\end{lstlisting}

This was a reproduced bypass of the selected v2.6.2 policy. It was not a
separate demonstration of UID~0 or target corruption, but it proved that the
timer could expire while the relevant socket and process remained active. The
control had therefore encoded elapsed time rather than the lifecycle of the
operation it was trying to constrain.

The first v2.7 lifecycle candidate was also retained rather than treated as a
successful fix. It recorded an \code{AEAD-CANDIDATE} and the later splice, but
the active state did not reach the permission decision; after the same
12-second delay, 32~bytes were still spliced. This intermediate failure showed
that merely adding socket-local state was insufficient unless the state
transition consulted by enforcement matched the AF\_ALG socket behaviour in
the test.

The frozen v2.7 source corrected the design in two ways:

\begin{enumerate}
    \item
    the AEAD association is counted for the lifetime of the tracked bound
    socket and released during socket teardown, with elapsed time retained only
    for event timestamps; and

    \item
    the splice flag is keyed by TID and exists only between
    \code{fentry/do\_splice} and \code{fexit/do\_splice} for that exact
    invocation.
\end{enumerate}

The retained regression again kept the AEAD operation socket open and waited
12~seconds before attempting the same SUID-to-pipe splice. This time the trace
contained \code{AEAD-ACTIVE}, \code{SPLICE-ARM}, and
\code{DENY-AEAD-SPLICE-SUID}; userspace received:

\begin{lstlisting}
v2.7 lifecycle regression: tgid=9531
AEAD socket active; waiting 12 seconds...
splice blocked: errno=1 (Operation not permitted)
\end{lstlisting}

The result addresses the specific expiry defect because the same delayed
sequence changed from allowed to denied after lifecycle state replaced the
timer. It does not prove that every socket-transfer, fork, exec, concurrency,
or teardown case is correct. Those cases were outside the preserved regression
matrix.

Cross-TGID isolation remained intentional. A v2.6.2 control held AEAD state in
one process while another process spliced the SUID target; the second process
successfully moved 32~bytes. The final v2.7 source retains process ownership for
AEAD state, so it does not create a host-wide correlation merely because two
unrelated processes perform the individual operations. This reduces broad
coupling between benign workloads, while also defining an evasion boundary if
an untested exploit implementation can distribute the required sequence across
processes.

\subsection{Adversarial Validation of the Frozen Candidate}

The final validation used the v2.7 executable whose digest begins
\code{7fe7b609}. The capture shows the userspace loader verifying the MORI
Monitor file, trusting TGID~702, attaching the enforcement engine, and
announcing that selective policy and structured telemetry were active. The
same executable remained running for the final Python and independent C tests;
the Phase~04 index records that it was not rebuilt between them.

\subsubsection{Python proof of concept}

The actual Python execution was performed after switching into the unprivileged
\code{labuser} account. A preliminary command in the same capture used an
incorrect home-directory expansion and failed before the test; it is not
counted as the security result. During the subsequent execution from
\code{/home/labuser}, the BPF trace recorded the following lifecycle for the
same TGID and TID:

\begin{lstlisting}
AEAD-ACTIVE
AEAD-ACCEPT-OBSERVED
SPLICE-ARM
DENY-AEAD-SPLICE-SUID
SPLICE-DISARM ret=-1
AEAD-RELEASE
\end{lstlisting}

The proof of concept was stopped by the LSM decision:

\begin{lstlisting}
PermissionError: [Errno 1] Operation not permitted
\end{lstlisting}

MORI userspace received the structured deny event for UID~1001 and delivered
the optional terminal notification. The security-relevant result is the LSM
\code{EPERM}; the message is presentation only.

Post-test validation recorded \code{labuser} as UID~1001 and returned the
known-good digest for \code{/usr/bin/su}:

\begin{lstlisting}
c74311fe5636b7d7f9a56239fa8adeeab12ba86fe7d41b91afa85bf9bbdae78b
\end{lstlisting}

The target therefore did not enter the changed mounted-view state used by MORI
Monitor in Chapter~6, and the test did not produce the UID~0 outcome seen during
the unmitigated reproduction.

\subsubsection{Independent C proof of concept}

The second validation used a separately implemented C proof of concept. Its
preserved provenance records the following identifiers:

\begin{lstlisting}
commit  ba3f503b1497350a2df3a89fab0427b4f868fe59
source  e1dec1348c0d43ab47d8f992d9a0336216fd37ff83e6fa1d9e39a80890fbd165
binary  3b2605dc5e820f40645f682385194469f9189cbb65ee99153b52aedf45be5706
\end{lstlisting}

Some provenance commands at the top of the combined validation capture were
run from the wrong directory and failed. They are not used for these values;
the separate provenance capture records the successful commit and hash output.

The C test reported that its first attempted write failed. The trace again
showed \code{AEAD-ACTIVE}, an observed accept, \code{SPLICE-ARM}, the v2.7
deny event, splice disarm with return value~$-1$, and AEAD release. After the
test, the account remained UID~1001 and \code{/usr/bin/su} retained the same
known-good digest. No privileged shell was observed.

\begin{table}[H]
    \centering
    \caption{Final adversarial results for the frozen MORI v2.7 candidate.}
    \label{tab:mori-v27-adversarial-results}
    \small
    \begin{tabularx}{\textwidth}{@{}>{\bfseries\raggedright\arraybackslash}p{0.20\textwidth}>{\raggedright\arraybackslash}p{0.24\textwidth}>{\raggedright\arraybackslash}p{0.23\textwidth}Y@{}}
        \toprule
        Test & Denial observation & Post-test identity and target & Conclusion \\
        \midrule
        Python PoC &
        Correlated splice returned \code{EPERM}; matching kernel and ring-buffer
        events were visible. &
        UID~1001; \code{/usr/bin/su} hash matched baseline. &
        Known Python path prevented. \\

        Independent C PoC &
        First attempted write failed after the matching v2.7 deny lifecycle. &
        UID~1001; \code{/usr/bin/su} hash matched baseline. &
        Second implementation path prevented. \\
        \bottomrule
    \end{tabularx}
\end{table}

The two implementations reduce dependence on one script's error handling or
layout. They do not make the sample size large and do not establish resistance
to all Copy Fail variants. The strongest supported statement is that the same
frozen v2.7 candidate interrupted both tested same-process paths before the
documented target change and privilege transition.

\subsection{Compatibility and Negative Controls}

Compatibility was evaluated incrementally rather than inferred from the final
denial alone. The most important contrast is between v1 and v2: v1 made even a
benign \code{gcm(aes)} bind unavailable, whereas the v2 series kept
\code{algif\_aead} loaded and allowed the tested benign AEAD operation before
selectively denying the correlated SUID splice.

The preserved controls are summarised in
\cref{tab:mori-compatibility-controls}. Version labels are retained because not
every earlier negative control was repeated in a separate capture against the
frozen v2.7 binary.

\begin{table}[H]
    \centering
    \caption{Compatibility, negative, and regression controls across MORI development.}
    \label{tab:mori-compatibility-controls}
    \scriptsize
    \begin{tabularx}{\textwidth}{@{}>{\bfseries}p{0.12\textwidth}>{\raggedright\arraybackslash}p{0.32\textwidth}>{\raggedright\arraybackslash}p{0.20\textwidth}Y@{}}
        \toprule
        Version & Control & Direct result & Meaning \\
        \midrule
        v1 & Unprivileged \code{gcm(aes)} AF\_ALG AEAD bind & Failed &
        Broad interface removal imposed a demonstrated compatibility cost. \\

        v2.2 & Same benign AEAD bind with v1 disabled & Succeeded and was
        observed &
        Observation did not itself require disabling AEAD. \\

        v2.5 & AEAD-only and SUID-splice-only shadow tests & Allowed; no
        \code{WOULD-DENY} &
        Neither individual signal was sufficient for the selected policy. \\

        v2.6 & Combined enforcement test & AEAD bind succeeded; correlated
        splice returned \code{EPERM} &
        The denial boundary could be placed after benign interface access. \\

        v2.6.2 & Cross-TGID AEAD and splice & Splice allowed &
        State was isolated by process rather than combined host-wide. \\

        v2.6.2 & SUID-target sweep & Attempts were denied; ten recorded target
        hashes remained unchanged &
        The predicate was not hard-coded only to \code{/usr/bin/su}; this was a
        v2.6.2 regression, not a universal target claim. \\

        v2.6.2 & Same-TGID operation after 12-second delay & Splice allowed &
        Reproduced a timer-based policy bypass. \\

        v2.7 & Same delayed lifecycle regression & Splice denied with
        \code{EPERM} &
        Socket-lifetime state addressed the reproduced expiry defect. \\

        v2.7 & Python and independent C PoCs & Both denied &
        The frozen candidate interrupted both final implementation paths. \\
        \bottomrule
    \end{tabularx}
\end{table}

The ten-target post-sweep capture lists \code{chfn}, \code{chsh},
\code{fusermount3}, \code{gpasswd}, \code{mount}, \code{newgrp},
\code{passwd}, \code{su}, \code{sudo}, and \code{umount} as unchanged. The
test demonstrates that the root-owner and SUID predicates could operate across
multiple recorded targets and that those targets retained their baseline
digests after the denied operations. It does not show that each program is an
independently exploitable Copy Fail target.

The benign AEAD evidence must also be interpreted narrowly. The v2.2 control
established a successful unprivileged \code{gcm(aes)} bind, and later controls
recorded bind and accept availability where stated. The dataset does not
preserve a complete encryption/decryption round trip or a representative
application workload. ``Benign functionality preserved'' therefore means that
MORI v2 did not remove the tested AF\_ALG AEAD interface operations and left
the isolated negative controls available. It does not mean that every legitimate
AEAD consumer was tested or that the control has zero compatibility impact.

\subsection{Operational and Security Limitations}

MORI v2.7 is an experimental, privileged, host-local compensating control. Its
result is useful because it is selective relative to v1, but several boundaries
prevent it from being treated as a general Linux exploit-prevention system.

\begin{itemize}
    \item
    \textbf{The vulnerable code remains present.} MORI denies an observed
    combination before the tested exploit reaches its target; it does not
    correct the \code{algif\_aead} implementation. Detaching the BPF programs,
    changing the loader, or booting without the required policy restores the
    underlying exposure.

    \item
    \textbf{The policy recognises behaviour, not intent.} AF\_ALG AEAD use and
    SUID-to-pipe splice activity can each be legitimate. A benign program that
    reproduces the complete same-process and same-thread combination may receive
    \code{EPERM}.

    \item
    \textbf{Process isolation creates an evasion boundary.} AEAD state is
    associated with an owner TGID. Cross-TGID activity is intentionally allowed.
    An untested implementation able to divide the required operations between
    processes, transfer descriptors, or otherwise change ownership semantics
    may fall outside the policy.

    \item
    \textbf{Bind observation precedes independent success confirmation.} The
    frozen source increments active state in \code{lsm/socket\_bind}. A later
    accept is separately observed, but enforcement consults the count created
    at bind. A failed bind whose socket remains open can therefore create
    conservative state until release. This was not evaluated across a corpus of
    failing benign applications.

    \item
    \textbf{Thread and lifecycle coverage is tested, not exhaustive.} TID-scoped
    splice state improves the precision of the permission decision, but fork,
    exec, descriptor passing, multithreaded races, abnormal teardown, and PID or
    TGID reuse were not subjected to a complete regression matrix.

    \item
    \textbf{Map-update failures allow the operation.} Source review shows that
    failure to allocate or update the active-AEAD or splice state is logged and
    the original operation is preserved. This avoids turning internal pressure
    into an uncontrolled system-wide denial, but it is fail-open for the
    compensating policy.

    \item
    \textbf{Monitor trust is local and start-up based.} The loader authenticates
    the monitor file and current service-cgroup processes before inserting TGIDs.
    The preserved source does not demonstrate continuous revalidation or removal
    of a trusted entry when the process lifecycle changes.

    \item
    \textbf{Telemetry is best-effort.} Ring-buffer exhaustion or userspace
    failure does not weaken a decision whose state was successfully created, but
    it can remove structured context and the optional TTY notification. Kernel
    trace output is also diagnostic rather than an independently protected audit
    trail.

    \item
    \textbf{Deployment depends on kernel facilities.} The tested design requires
    BPF LSM attachment, BTF/CO-RE information, the selected LSM and fentry hooks,
    and sufficient administrative privilege. The frozen provenance records
    clang~18.1.3 and libbpf~1.3.0, but no cross-distribution, cross-architecture,
    or production-kernel portability matrix was collected.

    \item
    \textbf{Performance and false-positive rates were not measured.} The
    experiment contains targeted positive and negative controls, not a sustained
    representative workload, controlled overhead benchmark, or statistical
    false-positive study.

    \item
    \textbf{The final evidence remains bounded.} Two proof-of-concept
    implementations, one principal target, an alternate target, and the stated
    regression suite improve confidence in the selected path. They do not prove
    protection against every algorithm, target, sequencing variation, or future
    exploitation technique.
\end{itemize}

The attacker-facing message deserves no separate security credit. It is useful
for demonstrating that the structured event reached userspace, but it also
reveals that a policy intervened and depends on access to an interactive
terminal. Removing it would not change the LSM result; losing it does not mean
the denial failed.

These limitations do not invalidate the experiment. They define what the
control achieved: a narrower temporary denial than complete interface removal,
with a documented policy, reproducible negative controls, a retained bypass,
and a regression-tested correction. They also explain why a vendor kernel fix
remains the preferred remediation.

\subsection{Answer to RQ4 and Transition to Vendor Comparison}

RQ4 asks whether a defensive control can interrupt the known exploitation path
before the kernel is patched and whether it can do so without unnecessarily
removing benign functionality. The experiments produce two different answers.

MORI Guard v1 answers the prevention part affirmatively but the compatibility
part negatively. On the vulnerable \code{6.8.0-116-generic} runtime, it kept
\code{algif\_aead} unavailable through the tested \code{modprobe} path, caused
the Python PoC to fail, preserved UID~1001, and kept \code{/usr/bin/su}
unchanged. The same restriction also caused the benign \code{gcm(aes)} bind to
fail. It interrupted the exploit by removing the complete tested interface.

The final MORI v2.7 result supports an affirmative but bounded answer to both
parts. The frozen candidate left AF\_ALG AEAD available in the exercised
controls and allowed individual or cross-process signals that did not satisfy
the policy. It denied the correlated SUID-to-pipe operation with \code{EPERM}
for both the Python and independent C proof-of-concept paths. After each test,
\code{labuser} remained UID~1001 and the selected privileged target retained
its known-good digest. The retained 12-second regression also shows that the
final lifecycle design corrected a specific bypass present in v2.6.2.

The qualifier is essential. MORI v2.7 preserves more benign functionality than
v1 because it does not remove the AEAD interface and does not deny every
individual prerequisite. It does not establish that all benign applications
remain unaffected, and process separation, state-allocation failure, lifecycle
edge cases, or a variant outside the selected correlation may evade it. The
appropriate conclusion is therefore:

\begin{lstlisting}
For the tested same-process Copy Fail paths in the Phase 04 laboratory,
MORI v2.7 interrupted exploitation before privilege escalation while
preserving the benign interface operations and negative controls that
were exercised. It is a selective compensating control, not a patch and
not proof of universal prevention.
\end{lstlisting}

This chapter deliberately stops short of answering RQ5. MORI Guard v1 and
Canonical's temporary mitigation both restrict \code{algif\_aead}, whereas
MORI v2.7 attempts a narrower behavioural decision inside the active interface.
Chapter~8 compares those approaches directly: what each blocks, what each
preserves, the trust and deployment assumptions each introduces, and why the
custom result does not replace vendor-supported remediation.